\section{Gabarito e resoluções comentadas — Gestão e Governança de TI}

\begin{solucao}{14.01}\textbf{CERTO.} EDM é governança; APO, BAI, DSS e MEA são gestão.\end{solucao}
\begin{solucao}{14.02}\textbf{ERRADO.} O enquadramento depende do objeto, nível decisório e responsabilidade; ``monitorar'' aparece nos dois níveis.\end{solucao}
\begin{solucao}{14.03}\textbf{CERTO.} COBIT 2019 prevê tailoring com base em fatores de desenho.\end{solucao}
\begin{solucao}{14.04}\textbf{ERRADO.} Incident management restaura serviço; problem management trata causas e recorrência.\end{solucao}
\begin{solucao}{14.05}\textbf{CERTO.} O princípio recomenda aproveitar o que funciona, não congelar o estado atual.\end{solucao}
\begin{solucao}{14.06}\textbf{CERTO.} Apetite a risco é decisão institucional de governança.\end{solucao}
\begin{solucao}{14.07}\textbf{CERTO.} Sem definição e fonte, o número não é reproduzível nem confiável para accountability.\end{solucao}
\begin{solucao}{14.08}\textbf{ERRADO.} Sequence flow fica dentro do pool; entre participantes usa-se message flow.\end{solucao}
\begin{solucao}{14.09}\textbf{CERTO.} Conformidade é um piso dentro de determinado escopo, não garantia de risco zero.\end{solucao}
\begin{solucao}{14.10}\textbf{CERTO.} Governança de dados define direitos de decisão e responsabilidades sobre o ativo.\end{solucao}

\begin{solucao}{14.11}\textbf{A.} Avaliar opções, dirigir a escolha e monitorar valor/risco caracteriza EDM.\end{solucao}
\begin{solucao}{14.12}\textbf{B.} Planejamento e organização de capacidades se alinham a APO.\end{solucao}
\begin{solucao}{14.13}\textbf{CERTO.} A cascata cria rastreabilidade entre necessidades e objetivos COBIT.\end{solucao}
\begin{solucao}{14.14}\textbf{B.} Problem management atua sobre causas e known errors.\end{solucao}
\begin{solucao}{14.15}\textbf{CERTO.} Mudanças de baixo risco podem ser pré-autorizadas/automatizadas; o foco é taxa de sucesso e risco.\end{solucao}
\begin{solucao}{14.16}\textbf{B.} Identificação, análise e avaliação formam risk assessment.\end{solucao}
\begin{solucao}{14.17}\textbf{B.} Lista de iniciativas sem owner, objetivo e métrica não traduz estratégia em execução governada.\end{solucao}
\begin{solucao}{14.18}\textbf{CERTO.} Meta isolada pode provocar fechamento artificial de tickets.\end{solucao}
\begin{solucao}{14.19}\textbf{B.} BSC conecta objetivos e medidas, adaptando-se à missão e ao valor público.\end{solucao}
\begin{solucao}{14.20}\textbf{B.} Comunicação entre pools é representada por message flow.\end{solucao}
\begin{solucao}{14.21}\textbf{CERTO.} Otimizar vem antes de automatizar; automação pode acelerar desperdício.\end{solucao}
\begin{solucao}{14.22}\textbf{A.} Tailoring ajusta governança e processo ao porte, risco e incerteza.\end{solucao}
\begin{solucao}{14.23}\textbf{CERTO.} Sucesso inclui valor e resultados, além de restrições de entrega.\end{solucao}
\begin{solucao}{14.24}\textbf{C.} LAI e LGPD coexistem; a decisão exige análise jurídica e de minimização.\end{solucao}
\begin{solucao}{14.25}\textbf{CERTO.} Glossário com owner e processo decisório reduz divergência semântica.\end{solucao}

\begin{solucao}{14.26}\textbf{C.} I e II corretas; III viola tailoring e priorização por fatores de desenho.\end{solucao}
\begin{solucao}{14.27}\textbf{CERTO.} O verbo não basta para distinguir gestão de governança.\end{solucao}
\begin{solucao}{14.28}\textbf{B.} A organização restaura, mas não aprende nem reduz recorrência por ausência de problem management efetivo.\end{solucao}
\begin{solucao}{14.29}\textbf{B.} Primeiro elimina-se/ajusta-se desperdício e controles redundantes; depois automatiza-se o fluxo útil.\end{solucao}
\begin{solucao}{14.30}\textbf{CERTO.} Regra de exclusão manipulável compromete comparabilidade e confiança do KPI.\end{solucao}
\begin{solucao}{14.31}\textbf{B.} Portfólio precisa equilibrar obrigação legal, risco, valor e restrições, não apenas ROI.\end{solucao}
\begin{solucao}{14.32}\textbf{B.} Cumprir prazo/custo não garante benefício; adoção e valor fazem parte da avaliação.\end{solucao}
\begin{solucao}{14.33}\textbf{CERTO.} Muitos KPIs mal definidos geram ruído, inconsistência e baixa auditabilidade.\end{solucao}
\begin{solucao}{14.34}\textbf{A.} Gateway paralelo convergente sincroniza caminhos paralelos.\end{solucao}
\begin{solucao}{14.35}\textbf{B.} Conformidade não autoriza risco acima do apetite; controles adicionais podem ser necessários.\end{solucao}
\begin{solucao}{14.36}\textbf{CERTO.} Responsabilidade é compartilhada; cliente mantém controles sob seu domínio.\end{solucao}
\begin{solucao}{14.37}\textbf{A.} Master/reference data management e governança tratam visão corporativa de entidades.\end{solucao}
\begin{solucao}{14.38}\textbf{A.} Aceitação de risco exige autoridade e critérios, não decisão técnica isolada.\end{solucao}
\begin{solucao}{14.39}\textbf{CERTO.} Scrum opera em equipe/produto; governança organizacional continua necessária.\end{solucao}
\begin{solucao}{14.40}\textbf{B.} Indicadores devem estar ligados a objetivos e hipóteses de valor.\end{solucao}

\begin{solucao}{14.41}\textbf{CERTO.} Execução pontual não demonstra alinhamento nem benefício sem conexão estratégica.\end{solucao}
\begin{solucao}{14.42}\textbf{B.} O incentivo atual promove gaming. Métrica balanceada reduz otimização de prazo às custas de escopo e valor.\end{solucao}
\begin{solucao}{14.43}\textbf{B.} COBIT não prescreve implementar tudo igualmente; fatores de desenho orientam foco.\end{solucao}
\begin{solucao}{14.44}\textbf{CERTO.} Aprovar investimento sem monitorar benefício/risco deixa incompleto o ciclo Evaluate--Direct--Monitor.\end{solucao}
\begin{solucao}{14.45}\textbf{B.} Melhorar processo exige entender propósito e risco antes de eliminar ou automatizar controles.\end{solucao}
\begin{solucao}{14.46}\textbf{CERTO.} Indicador pode cumprir meta estatística e falhar em representar a experiência crítica.\end{solucao}
\begin{solucao}{14.47}\textbf{B.} O problema é semântico e de lineage; terceiro dashboard não resolve definição divergente.\end{solucao}
\begin{solucao}{14.48}\textbf{A.} Risco de serviço crítico precisa de decisão no nível de governança compatível com apetite e impacto.\end{solucao}
\begin{solucao}{14.49}\textbf{CERTO.} Transparência e proteção de dados devem ser conciliadas conforme o regime jurídico.\end{solucao}
\begin{solucao}{14.50}\textbf{C.} O achado liga condição, causa de governança, efeito na confiabilidade e ações concretas.\end{solucao}
